% ============================================================
%  cap04_multipli_divisori.tex — Matemagica
%  Capitolo 4: Multipli, divisori, MCD e mcm
% ============================================================

\chapter{Multipli, divisori, MCD e mcm}

% --- Obiettivi di apprendimento ---
\section*{Obiettivi di apprendimento}
\begin{itemize}[leftmargin=1.5em]
    \item Comprendere il concetto di multiplo e di divisore di un numero naturale.
    \item Rappresentare insiemi di multipli e di divisori con la notazione insiemistica.
    \item Riconoscere i numeri primi e distinguerli dai numeri composti.
    \item Scomporre un numero naturale in fattori primi.
    \item Applicare i criteri di divisibilità per 2, 3, 4, 5, 6, 9, 10 e 25.
    \item Calcolare il minimo comune multiplo (mcm) e il massimo comune divisore (MCD)
          di due o più numeri.
    \item Conoscere le principali proprietà di mcm e MCD e applicarle.
\end{itemize}

% ============================================================
\section{I multipli di un numero naturale}
% ============================================================

\begin{definizione}
    I \textbf{multipli} di un numero naturale $n$ sono tutti i numeri naturali che si
    ottengono moltiplicando $n$ per ciascun numero naturale:
    \[
        n \times 0,\quad n \times 1,\quad n \times 2,\quad n \times 3,\quad \ldots
    \]
    I multipli di $n$ formano un insieme, indicato con il simbolo $M_n$, la cui
    caratteristica è \emph{«essere multiplo di $n$»}.
\end{definizione}

\begin{esempio}
    Alcuni insiemi di multipli:
    \begin{align*}
        M_3    & = \{0;\; 3;\; 6;\; 9;\; 12;\; 15;\; 18;\; 21;\; \ldots\} \\
        M_5    & = \{0;\; 5;\; 10;\; 15;\; 20;\; 25;\; 30;\; \ldots\}     \\
        M_{13} & = \{0;\; 13;\; 26;\; 39;\; \ldots\}
    \end{align*}
\end{esempio}

\begin{attenzione}
    \textbf{Attenzione.} Lo zero è multiplo di qualsiasi numero naturale, perché
    $n \times 0 = 0$ per ogni $n$. Tuttavia, quando si cerca il minimo comune multiplo
    si considera il più piccolo multiplo \emph{diverso da zero}.
\end{attenzione}

\begin{esempio}
    Rappresentare per elencazione i seguenti insiemi, usando la notazione insiemistica:

    \medskip
    \textbf{A} $= \{x \in \mathbb{N} \mid x \in M_3 \text{ e } x > 24\}$

    Si cercano i multipli di 3 maggiori di 24:
    $27 = 3 \times 9$, $30 = 3 \times 10$, $33 = 3 \times 11$, \ldots

    \[
        A = \{27;\; 30;\; 33;\; 36;\; \ldots\}
    \]

    \medskip
    \textbf{B} $= \{x \in \mathbb{N} \mid x \in M_5 \text{ e } 10 < x \leq 30\}$

    Si cercano i multipli di 5 compresi tra 10 (escluso) e 30 (incluso):

    \[
        B = \{15;\; 20;\; 25;\; 30\}
    \]

    \medskip
    \textbf{C} $= \{x \in \mathbb{N} \mid x \in M_{13} \text{ e } x < 13\}$

    Il più piccolo multiplo di 13 diverso da zero è 13 stesso. Non esistono multipli
    di 13 (positivi) inferiori a 13, quindi:

    \[
        C = \emptyset
    \]
\end{esempio}

% ============================================================
\section{Il minimo comune multiplo (mcm)}
% ============================================================

\begin{esempio}
    \textbf{Problema introduttivo — I pianeti in fila.}

    Giove e Mercurio sono due pianeti del sistema solare. Come la Terra, si muovono
    ruotando attorno al Sole: Giove compie un'orbita completa in 12 anni, Mercurio
    in 88 giorni. Adattiamo il problema ai nostri strumenti: supponiamo che Giove
    impieghi \textbf{4 anni} e Mercurio \textbf{6 anni} per compiere un'orbita.
    Se quest'anno si trovassero allineati rispetto al Sole, dopo quanti anni si
    ripeterebbe il fenomeno?

    \medskip
    I momenti di allineamento sono i multipli \emph{comuni} alle due durate:
    \begin{align*}
        M_4 & = \{0;\; 4;\; 8;\; \mathbf{12};\; 16;\; \ldots\} \\
        M_6 & = \{0;\; 6;\; \mathbf{12};\; 18;\; \ldots\}
    \end{align*}
    Il primo multiplo comune diverso da zero è \textbf{12}: i due pianeti si
    riallineano dopo 12 anni.
\end{esempio}

\begin{definizione}
    Il più piccolo multiplo \emph{diverso da zero} comune a due o più numeri si
    chiama \textbf{minimo comune multiplo} e si indica con \textbf{mcm}.
\end{definizione}

\begin{esempio}
    Trovare il mcm di 3 e 7 per enumerazione:
    \begin{align*}
        M_7 & = \{0;\; 7;\; 14;\; \mathbf{21};\; 28;\; \ldots\}                 \\
        M_3 & = \{0;\; 3;\; 6;\; 9;\; 12;\; 15;\; 18;\; \mathbf{21};\; \ldots\}
    \end{align*}
    Fra tutti i multipli comuni a 7 e 3, il più piccolo diverso da zero è 21,
    quindi $\text{mcm}(3;\, 7) = 21$.
\end{esempio}

% ============================================================
\section{I divisori di un numero naturale}
% ============================================================

\begin{definizione}
    I \textbf{divisori} di un numero naturale $n$ sono tutti i numeri naturali che
    dividono $n$ esattamente, cioè con resto zero. I divisori di $n$ formano un
    insieme, indicato con il simbolo $D_n$, la cui caratteristica è
    \emph{«essere divisore di $n$»}.
\end{definizione}

\begin{esempio}
    Alcuni insiemi di divisori:
    \begin{align*}
        D_6    & = \{1;\; 2;\; 3;\; 6\}                        \\
        D_{11} & = \{1;\; 11\}                                 \\
        D_{40} & = \{1;\; 2;\; 4;\; 5;\; 8;\; 10;\; 20;\; 40\}
    \end{align*}
\end{esempio}

\begin{esempio}
    Rappresentare per elencazione i seguenti insiemi:

    \medskip
    \textbf{A} $= \{x \in \mathbb{N} \mid x \in D_{15}\}$

    I divisori di 15 sono: $1 \times 15 = 15$, $3 \times 5 = 15$.
    \[
        A = \{1;\; 3;\; 5;\; 15\}
    \]

    \medskip
    \textbf{B} $= \{x \in \mathbb{N} \mid x \in D_{31}\}$

    Il numero 31 è divisibile solo per 1 e per se stesso (come vedremo, è un numero primo).
    \[
        B = \{1;\; 31\}
    \]

    \medskip
    \textbf{C} $= \{x \in \mathbb{N} \mid x \in D_{46} \text{ e } x < 23\}$

    I divisori di 46 sono: $\{1;\; 2;\; 23;\; 46\}$. Fra questi, quelli minori di 23 sono:
    \[
        C = \{1;\; 2\}
    \]
\end{esempio}

\begin{attenzione}
    \textbf{I divisori sono in numero finito.} Ogni numero naturale $n > 0$ ha un
    numero limitato di divisori: essi sono compresi tra 1 e $n$.
    I multipli di un numero, invece, sono infiniti.
\end{attenzione}

% ============================================================
\section{Numeri primi e numeri composti}
% ============================================================

Osservando gli insiemi dei divisori, emerge una distinzione importante.

\begin{definizione}
    Un numero naturale maggiore di 1 si dice \textbf{numero primo} se ha esattamente
    due divisori: 1 e se stesso.

    Un numero naturale maggiore di 1 che ha più di due divisori si dice
    \textbf{numero composto}.

    Il numero 1 non è né primo né composto: è un caso a parte, con un solo divisore.
\end{definizione}

\begin{esempio}
    \begin{itemize}[leftmargin=1.5em]
        \item $D_7 = \{1;\; 7\}$ \quad — 7 è \textbf{primo} \quad\textit{(corretto)}
        \item $D_{13} = \{1;\; 13\}$ \quad — 13 è \textbf{primo} \quad\textit{(corretto)}
        \item $D_{12} = \{1;\; 2;\; 3;\; 4;\; 6;\; 12\}$ \quad — 12 è \textbf{composto} \quad\textit{(corretto)}
        \item $D_1 = \{1\}$ \quad — 1 non è né primo né composto
    \end{itemize}
\end{esempio}

I primi numeri primi sono:
\[
    2,\; 3,\; 5,\; 7,\; 11,\; 13,\; 17,\; 19,\; 23,\; 29,\; 31,\; 37,\; 41,\; 43,\; 47,\; \ldots
\]

\begin{attenzione}
    \textbf{Il 2 è l'unico numero primo pari.} Ogni altro numero pari è divisibile per 2
    (oltre che per 1 e per se stesso), quindi ha più di due divisori ed è composto.
\end{attenzione}

\begin{sapevi}
    \textbf{Il crivello di Eratostene.}
    Eratostene di Cirene (276–194 a.C.), matematico e geografo greco, ideò un metodo
    elegante per trovare tutti i numeri primi fino a un certo valore. Si scrivono
    tutti i numeri da 2 in poi; si tiene il 2 (primo) e si eliminano tutti i suoi
    multipli; si tiene il successivo non eliminato (3) e si eliminano tutti i suoi
    multipli; e così via. I numeri che rimangono sono esattamente i numeri primi.
    Questo metodo è ancora oggi usato in informatica per generare liste di numeri primi.
\end{sapevi}

% ============================================================
\section{La scomposizione in fattori primi}
% ============================================================

\begin{definizione}
    Ogni numero naturale maggiore di 1 può essere scritto come prodotto di numeri
    primi, detti \textbf{fattori primi}. Questa scrittura si chiama
    \textbf{scomposizione in fattori primi} (o \textbf{fattorizzazione}).
    La scomposizione in fattori primi di un numero è \emph{unica}
    (a meno dell'ordine dei fattori).
\end{definizione}

Il metodo pratico consiste nel dividere il numero per il più piccolo fattore
primo che lo divide esattamente, poi ripetere con il quoziente ottenuto, fino
ad arrivare a 1.

\begin{esempio}
    Scomporre 60 in fattori primi.

    Si procede con divisioni successive:

    \[
        \begin{array}{r|l}
            60 & 2 \\
            30 & 2 \\
            15 & 3 \\
            5  & 5 \\
            1  &
        \end{array}
    \]

    Quindi: $60 = 2 \times 2 \times 3 \times 5 = 2^2 \times 3 \times 5$
    \quad\textit{(corretto)}
\end{esempio}

\begin{esempio}
    Scomporre 84 in fattori primi.

    \[
        \begin{array}{r|l}
            84 & 2 \\
            42 & 2 \\
            21 & 3 \\
            7  & 7 \\
            1  &
        \end{array}
    \]

    Quindi: $84 = 2^2 \times 3 \times 7$
    \quad\textit{(corretto)}
\end{esempio}

\begin{attenzione}
    \textbf{Come scegliere il divisore primo.} Si comincia sempre dal numero primo
    più piccolo (2), e si passa al successivo solo quando il quoziente non è più
    divisibile per il primo corrente. L'ordine è: 2, 3, 5, 7, 11, 13, \ldots

    Se non si è sicuri che un numero sia primo, si prova a dividerlo per tutti i
    numeri primi fino alla sua radice quadrata: se nessuno lo divide, allora è primo.
\end{attenzione}

% ============================================================
\section{I criteri di divisibilità}
% ============================================================

Trovare i divisori di un numero grande può essere complicato. I \textbf{criteri
    di divisibilità} permettono di stabilire rapidamente se un numero è divisibile
per un altro, senza eseguire la divisione.

\begin{regola}
    \begin{itemize}[leftmargin=1.5em]
        \item Un numero è divisibile per \textbf{2} se è pari (cioè termina con
              0, 2, 4, 6 o 8).
        \item Un numero è divisibile per \textbf{3} se la somma delle sue cifre
              è divisibile per 3.
        \item Un numero è divisibile per \textbf{4} se il numero formato dalle
              sue ultime due cifre è divisibile per 4.
        \item Un numero è divisibile per \textbf{5} se termina con 0 o con 5.
        \item Un numero è divisibile per \textbf{6} se è divisibile per 2 e per 3
              contemporaneamente.
        \item Un numero è divisibile per \textbf{9} se la somma delle sue cifre
              è divisibile per 9.
        \item Un numero è divisibile per \textbf{10} se termina con 0.
        \item Un numero è divisibile per \textbf{25} se termina con 00, 25, 50 o 75.
    \end{itemize}
\end{regola}

\begin{esempio}
    Verificare la divisibilità di 5\,748 applicando i criteri.

    \begin{itemize}[leftmargin=1.5em]
        \item \textbf{Per 2:} termina con 8 (cifra pari) \quad$\Rightarrow$\quad
              divisibile \quad\textit{(corretto)}
        \item \textbf{Per 3:} $5 + 7 + 4 + 8 = 24$, e $24 \div 3 = 8$ \quad$\Rightarrow$\quad
              divisibile \quad\textit{(corretto)}
        \item \textbf{Per 4:} le ultime due cifre sono 48, e $48 \div 4 = 12$ \quad$\Rightarrow$\quad
              divisibile \quad\textit{(corretto)}
        \item \textbf{Per 5:} termina con 8 \quad$\Rightarrow$\quad
              non divisibile
        \item \textbf{Per 6:} divisibile per 2 e per 3 \quad$\Rightarrow$\quad
              divisibile \quad\textit{(corretto)}
        \item \textbf{Per 9:} $5 + 7 + 4 + 8 = 24$, e $24 \div 9$ non è esatto \quad$\Rightarrow$\quad
              non divisibile
        \item \textbf{Per 10:} non termina con 0 \quad$\Rightarrow$\quad
              non divisibile
    \end{itemize}
\end{esempio}

\begin{sapevi}
    \textbf{Perché funziona il criterio del 9?}
    Ogni numero può essere scritto come somma di potenze di 10:
    per esempio $345 = 3 \times 100 + 4 \times 10 + 5$. Poiché
    $10 = 9 + 1$, $100 = 99 + 1$, ecc., si può dimostrare che un numero e la
    somma delle sue cifre danno lo stesso resto quando vengono divisi per 9.
    Lo stesso ragionamento vale per il criterio del 3.
\end{sapevi}

% ============================================================
\section{Il massimo comune divisore (MCD)}
% ============================================================

\begin{definizione}
    Il più grande divisore comune a due o più numeri si chiama
    \textbf{massimo comune divisore} e si indica con \textbf{MCD}.
\end{definizione}

\begin{esempio}
    Trovare il MCD di 16 e 64 per enumerazione:
    \begin{align*}
        D_{16} & = \{1;\; 2;\; 4;\; 8;\; \mathbf{16}\}             \\
        D_{64} & = \{1;\; 2;\; 4;\; 8;\; \mathbf{16};\; 32;\; 64\}
    \end{align*}
    Fra tutti i divisori comuni a 16 e 64, il più grande è 16,
    quindi $\text{MCD}(16;\, 64) = 16$.
\end{esempio}

% ============================================================
\section{Calcolo di mcm e MCD con la scomposizione in fattori primi}
% ============================================================

Il metodo per enumerazione funziona bene con numeri piccoli, ma diventa
impraticabile con numeri grandi. Il metodo più efficiente si basa sulla
scomposizione in fattori primi.

\begin{regola}
    \textbf{Calcolo del MCD.}
    \begin{enumerate}[leftmargin=1.5em]
        \item Scomporre ciascun numero in fattori primi.
        \item Individuare i fattori primi \emph{comuni} a tutte le scomposizioni.
        \item Il MCD è il prodotto di quei fattori comuni, ciascuno preso con
              l'\emph{esponente minore}.
    \end{enumerate}

    \textbf{Calcolo del mcm.}
    \begin{enumerate}[leftmargin=1.5em]
        \item Scomporre ciascun numero in fattori primi.
        \item Raccogliere \emph{tutti} i fattori primi che compaiono in almeno
              una scomposizione.
        \item Il mcm è il prodotto di tutti quei fattori, ciascuno preso con
              l'\emph{esponente maggiore}.
    \end{enumerate}
\end{regola}

\begin{esempio}
    Calcolare MCD e mcm di 180 e 252.

    \medskip
    \textbf{Passo 1 — scomposizione in fattori primi:}

    \[
        \begin{array}{r|l}
            180 & 2 \\
            90  & 2 \\
            45  & 3 \\
            15  & 3 \\
            5   & 5 \\
            1   &
        \end{array}
        \qquad\Rightarrow\qquad
        180 = 2^2 \times 3^2 \times 5
    \]

    \[
        \begin{array}{r|l}
            252 & 2 \\
            126 & 2 \\
            63  & 3 \\
            21  & 3 \\
            7   & 7 \\
            1   &
        \end{array}
        \qquad\Rightarrow\qquad
        252 = 2^2 \times 3^2 \times 7
    \]

    \medskip
    \textbf{Passo 2 — MCD} (fattori comuni con esponente minore):

    I fattori comuni sono 2 e 3. L'esponente minore di 2 è 2; l'esponente minore di 3 è 2.

    \[
        \text{MCD}(180;\, 252) = 2^2 \times 3^2 = 4 \times 9 = 36
    \]

    \textbf{Passo 3 — mcm} (tutti i fattori con esponente maggiore):

    I fattori che compaiono sono 2, 3, 5 e 7.

    \[
        \text{mcm}(180;\, 252) = 2^2 \times 3^2 \times 5 \times 7 = 4 \times 9 \times 5 \times 7 = 1260
    \]
\end{esempio}

\begin{attenzione}
    \textbf{Errore comune.} Nel calcolo del MCD si prendono \emph{solo} i fattori
    \emph{comuni} con l'esponente \emph{minore}. Nel calcolo del mcm si prendono
    \emph{tutti} i fattori (anche quelli non comuni) con l'esponente \emph{maggiore}.
    Confondere i due criteri è l'errore più frequente.
\end{attenzione}

% ============================================================
\section{Proprietà di mcm e MCD}
% ============================================================

\begin{regola}
    Le seguenti proprietà permettono di semplificare molti calcoli:

    \begin{enumerate}[leftmargin=1.5em]
        \item \textbf{Numeri coprimi.} Se MCD$(a;\, b) = 1$, i due numeri si dicono
              \textbf{coprimi} (o primi tra loro). In questo caso:
              \[
                  \text{mcm}(a;\, b) = a \times b
              \]
              \textit{Esempio:} MCD$(4;\, 9) = 1$, quindi mcm$(4;\, 9) = 36$.

        \item \textbf{Se $a$ divide $b$.} Se $a$ è divisore di $b$, allora:
              \[
                  \text{MCD}(a;\, b) = a \qquad \text{e} \qquad \text{mcm}(a;\, b) = b
              \]
              \textit{Esempio:} MCD$(4;\, 12) = 4$ e mcm$(4;\, 12) = 12$.

        \item \textbf{Relazione tra MCD e mcm.} Per due numeri qualsiasi $a$ e $b$:
              \[
                  \text{MCD}(a;\, b) \times \text{mcm}(a;\, b) = a \times b
              \]
              \textit{Esempio:} $a = 12$, $b = 18$:
              $\text{MCD} = 6$, $\text{mcm} = 36$,
              e infatti $6 \times 36 = 216 = 12 \times 18$. \quad\textit{(corretto)}

        \item \textbf{Il MCD è sempre minore o uguale al più piccolo dei numeri.}
              Il mcm è sempre maggiore o uguale al più grande dei numeri.
    \end{enumerate}
\end{regola}

\begin{sapevi}
    \textbf{mcm e MCD nella vita quotidiana.}
    Il minimo comune multiplo non serve solo per i pianeti. Si usa ogni volta che
    bisogna sincronizzare eventi che si ripetono con periodicità diversa: i turni
    di lavoro, le finestre di manutenzione di macchine industriali, oppure —
    in musica — per trovare il momento in cui due ritmi diversi tornano a coincidere.
    Il MCD, invece, è alla base del calcolo delle frazioni equivalenti e della
    semplificazione delle frazioni, concetti che affronteremo nel prossimo capitolo.
\end{sapevi}

% ============================================================
\section*{Riepilogo}
% ============================================================

\begin{regola}
    \textbf{Riepilogo del capitolo.}

    \begin{itemize}[leftmargin=1.5em]
        \item I \textbf{multipli} di $n$ sono i prodotti $n \times 0,\; n \times 1,\;
                  n \times 2,\; \ldots$ e formano l'insieme $M_n$ (infinito).
        \item I \textbf{divisori} di $n$ sono i numeri naturali che dividono $n$
              esattamente e formano l'insieme $D_n$ (finito).
        \item Un numero con esattamente due divisori è \textbf{primo}; con più di
              due è \textbf{composto}; il numero 1 non è né primo né composto.
        \item Ogni numero composto si scompone in modo unico come prodotto di
              \textbf{fattori primi}.
        \item I \textbf{criteri di divisibilità} permettono di riconoscere
              rapidamente i divisori di 2, 3, 4, 5, 6, 9, 10 e 25.
        \item Il \textbf{MCD} si calcola moltiplicando i fattori primi comuni con
              l'esponente \emph{minore}.
        \item Il \textbf{mcm} si calcola moltiplicando tutti i fattori primi con
              l'esponente \emph{maggiore}.
        \item Per due numeri $a$ e $b$: $\text{MCD}(a;\,b) \times \text{mcm}(a;\,b)
                  = a \times b$.
    \end{itemize}
\end{regola}